\setcounter{section}{5}
\section{拟阵与贪心结构}
\sourcepage{1830}
本节和下一节讨论拟阵及两类相关结构——有向拟阵和贪心结构——的某些拓扑方面。基本定义见 Welsh 的第9章。关于拟阵复形及几何格的其他拓扑事实见(11.10)，另见 Björner (1992)。
\subsection*{基复形与图的划分}
以下结果由 E. Győry 和 L. Lovász 回答 A. Frank 和 S. Maurer 的猜想时证明。
\entry{定理6.1}{（Lovász 1977，Győry 1978）}设 $G=(V,E)$ 是 $k$ 连通图，$\{v_1,\ldots,v_k\}$ 是 $k$ 个顶点，$n_1,\ldots,n_k$ 是正整数且 $\sum_i n_i=|V|$。则存在 $V$ 的划分 $\{V_1,\ldots,V_k\}$，使对每个 $i=1,\ldots,k$，$v_i\in V_i$、$|V_i|=n_i$，且 $V_i$ 在 $G$ 中诱导连通子图。
Lovász 的证明使用拓扑方法，Győry 的证明则没有。本节末将概要介绍 $k=3$ 的 Lovász 证明，以说明拓扑推理的用法。证明依赖于与 $G$ 中某些森林相关的多面体复形的连通性。类似复形可在拟阵的基上、乃至更一般地在贪心结构的基上定义。贪心结构的表述包含前两者，我们用它建立一般结果。先回顾定义。
集合系统 $(E,\mathcal F)$，其中 $\mathcal F\subseteq2^E$，若满足以下公理，称为\textbf{贪心结构}：
\begin{description}
\item[(G1)] $\varnothing\in\mathcal F$；
\item[(G2)] 对每个非空 $A\in\mathcal F$，存在 $x\in A$，使 $A- x\in\mathcal F$；
\item[(G3)] 若 $A,B\in\mathcal F$ 且 $|A|>|B|$，则存在 $x\in A-B$，使 $B\cup x\in\mathcal F$。
\end{description}
若还满足以下条件，则称为\textbf{区间贪心结构}：
\begin{description}\item[(G4)] 若 $A\subseteq B\subseteq C$，其中 $A,B,C\in\mathcal F$，且对某个 $x\in E-C$ 有 $A\cup x,C\cup x\in\mathcal F$，则 $B\cup x\in\mathcal F$。\end{description}
$\mathcal F$ 中的集合称为可行集，极大可行集称为基。所有基的基数相同，记为 $r$，称为贪心结构的秩。
本文仅涉及拟阵（可行集为独立集）和有根图的分枝贪心结构（可行集为构成包含根的树的边集）。二者都是区间贪心结构。其他例子及更多内容见 Welsh 的第9章以及 Korte、Lovász 和 Schrader (1991)、Björner 和 Ziegler (1992) 的综述。
\sourcepage{1831}
除拟阵情形外，贪心结构的可行集不构成单纯复形。不过，按包含关系排序的偏序集 $\overline{\mathcal F}=\mathcal F-\{\varnothing\}$ 的序复形给出了有用的拓扑。称秩为 $r$ 的贪心结构 $(E,\mathcal F)$ 为 $k$ 连通的，如果对每个 $A\in\mathcal F$，存在 $B\in\mathcal F$，使 $A\subseteq B$、$|B-A|=\min(k,r-|A|)$，且每个满足 $A\subseteq C\subseteq B$ 的 $C$ 都属于 $\mathcal F$。拟阵是 $r$ 连通的；$k$ 连通有根图的分枝贪心结构也是 $k$ 连通的。
\entry{命题6.2}{（Björner、Korte 和 Lovász 1985）}设 $(E,\mathcal F)$ 是 $k$ 连通区间贪心结构（$k\ge2$），则其可行集偏序集 $(\overline{\mathcal F},\subseteq)$ 是拓扑意义下的 $(k-2)$ 连通的。
此结果通过定理10.8由(11.10)(iii)推出：$\overline{\mathcal F}$ 的极小元素组成的截集 $C$ 所对应的截集复形 $\Gamma(\overline{\mathcal F},C)$ 是秩至少为 $k$ 的拟阵复形。
设 $\mathcal B$ 为秩 $r$ 的贪心结构的所有基。若 $B_1\cap B_2\in\mathcal F$ 且 $|B_1\cap B_2|=r-1$，称两基相邻。连接所有相邻基，得到以 $\mathcal B$ 为顶点集的\textbf{基图}。基图中最短的圈可明确描述，有两类三角形和一类四边形：
\entry{6.3}{}三个基 $A\cup x,A\cup y,A\cup z$，其中 $A\in\mathcal F$、$|A|=r-1$，张成第一类三角形。
\entry{6.4}{}三个基 $A\cup x\cup y,A\cup x\cup z,A\cup y\cup z$，其中 $A\in\mathcal F$、$|A|=r-2$，张成第二类三角形。
\entry{6.5}{}四个基 $A\cup x\cup u,A\cup x\cup v,A\cup y\cup u,A\cup y\cup v$，其中 $A\in\mathcal F$、$|A|=r-2$，张成一个四边形。
分枝贪心结构不会出现第二类三角形。现在在每个三角形及四边形内贴上一个二维胞腔（“膜”），得到二维正则胞腔复形 $\mathcal K$，称为\textbf{基复形}。不难检验，秩不超过2的任何2连通贪心结构，其基复形都是1连通的，即连通且单连通。秩为2时，这直接由交换公理(G3)得到。较高秩时有：
\entry{定理6.6}{（Björner、Korte 和 Lovász 1985）}任何3连通区间贪心结构的基复形 $\mathcal K$ 都是1连通的。
为了说明第二部分的一些工具，给出简短证明。设 $P$ 是 $\mathcal K$ 的闭胞腔按包含关系组成的偏序集，$Q$ 是 $(\mathcal F,\subseteq)$ 的最高三层，即秩分别为 $r-2,r-1,r$ 的可行集。令 $f:P\to Q$ 为反序映射，把每个胞腔映到张成它的各个基的交。由命题6.2和引理11.12，$Q$ 是1连通的；由定理10.5，只需检验所有纤维 $f^{-1}(Q_{\ge A})$ 是1连通的。
\sourcepage{1832}
若 $r(A)=r-i$，$i=0,1,2$，则 $f^{-1}(Q_{\ge A})$ 是收缩 $A$ 所得秩为 $i$ 的贪心结构的基复形；我们已检验秩不超过2的情形是1连通的。
设 $P=B_1B_2\cdots B_d$、$Q=B_dB_{d+1}\cdots B_g$ 是拟阵基图中的路，令 $PQ=B_1\cdots B_g$ 为其拼接。若 $R$ 形如 $BCB$、$BCDB$ 或 $BCDEB$，其中 $B=B_d$，则称 $PQ$ 与 $PRQ$ 相差一个\textbf{初等同伦}。
\entry{定理6.7}{（Maurer 1973）}拟阵基图中具有相同端点的任意两条路 $P,P'$，可通过一系列初等同伦相互变换。
Maurer 的同伦定理显然是定理6.6在拟阵情形下的组合表述。下一节将给出其在有向拟阵中的应用。
现在回到定理6.1。下面引用 Lovász (1979) 对 $k=3$ 的证明概要，方括号中作了若干调整以适应本文记号。
设 $G$ 是3连通图，$v_1,v_2,v_3\in V(G)$，且 $n_1+n_2+n_3=|V(G)|$。加入新点 $a$ 并将其与 $v_1,v_2,v_3$ 相连。考虑为新图 $G'$ 构造的拓扑空间 $\mathcal K$。［按本文语言，$\mathcal K$ 是有根图 $(G',a)$ 所确定的分枝贪心结构的基复形；这个贪心结构是3连通的，其基是 $G'$ 的生成树。］对每棵生成树 $T$，令 $f_i(T)$ 为从 $a$ 沿边 $(a,v_i)$ 在 $T$ 中可达的点数（$i=1,2$）。映射 $f:T\mapsto(f_1(T),f_2(T))$ 将 $\mathcal K$ 的顶点映为平面格点。对每个四边形二维胞腔加一条对角线，将其分成两个三角形，得到三角剖分 $\widetilde{\mathcal K}$。在每个三角形上线性延拓 $f$，得到到平面的连续映射。其像显然包含于三角形 $\Delta=\{x\ge0,y\ge0,x+y\le n\}$；我们证明该映射覆盖 $\Delta$。
先取生成树 $T_1,T_2,T_3$，使 $f(T_1)=(n,0)$、$f(T_2)=(0,n)$、$f(T_3)=(0,0)$。应用2连通贪心结构的基图连通这一事实，在 $G'-(a,v_3)$ 中选取连接 $T_1,T_2$ 的折线 $P_{12}$，使其所有点满足 $f_3=0$。于是 $f(P_{12})$ 沿 $\Delta$ 的边连接 $(n,0)$ 与 $(0,n)$。类似定义 $P_{23},P_{31}$。
由定理6.6，$P_{12}+P_{23}+P_{31}$ 可在 $\mathcal K$ 中收缩为一点。因此其像也可在 $f(\mathcal K)$ 中收缩为一点。但三角形边界不能在去掉一个内部点的三角形中收缩为一点（Brouwer 定理13.1），故 $f(\mathcal K)$ 必覆盖整个 $\Delta$。特别地，$(n_1,n_2)$ 属于某个三角形的像。容易看出，这意味着它属于某个顶点的像：即存在生成树 $T$ 满足 $f_1(T)=n_1$、$f_2(T)=n_2$。
\sourcepage{1833}
于是 $T-a$ 的三个分支给出所需的 $V(G)$ 的划分。
定理6.6是更一般结果的特例：对任意 $k$ 连通区间贪心结构，某个高维基复形是 $(k-2)$ 连通的。将上述 $k=3$ 的思想推广，即得一般的定理6.1。详情见 Lovász (1977) 和 Björner、Korte、Lovász (1985)。
\subsection*{Tutte 同伦定理}
若拟阵能在每个域上坐标化，称其为\textbf{正则拟阵}。Tutte (1958) 用排除子式刻画了正则拟阵，其证明本质上依赖一个“同伦定理”：某些二维复形是1连通的。R. Reid 和 R. Bixby 也使用它证明了在 $GF(3)$ 上可表示性的排除子式刻画。后来 P. Seymour 等人找到了不使用此定理的其他证明；见 Seymour 的第10章。
Tutte 同伦定理似乎是组合学中这类拓扑结果最早的例子。完整叙述及证明都相当繁琐，这里只陈述足以理解其性质的内容；详情见 Tutte (1958, 1965)。
设 $L$ 是秩 $r$ 的有限几何格，$L^i$ 是秩为 $i$ 的平坦集之集；$L^{r-1},L^{r-2},L^{r-3}$ 的元素分别称为余点、余线、余平面。把平坦集 $X$ 视为点集 $L^1$ 的子集 $\bar X=\{p\in L^1:p\le X\}$。给定点 $a\in L^1$，在顶点集 $L^{r-1}_{\not\ge a}=\{X\in L^{r-1}:X\not\ge a\}$ 上定义图 $TG(L,a)$：两个不含 $a$ 的余点 $X,Y$ 相邻，当且仅当 $X\wedge Y$ 是余线且 $\bar X\cup\bar Y\ne L^1-a$。在该图的下列三角形和四边形中贴上二维胞腔，构造二维正则胞腔复形 $TC(L,a)$：
\entry{6.8}{}满足 $\rk(X\wedge Y\wedge Z)\ge r-3$ 的三角形 $XYZX$。
\entry{6.9}{}四边形 $XYZTX$，其中 $P=X\wedge Y\wedge Z\wedge T$ 满足 $\rk(P)=r-3$，且余线 $P\vee a$ 恰被两个余点覆盖，或者区间 $[P,\hat1]$ 同构于从 Fano 拟阵 $F_7$ 中删去一个点所得拟阵的平坦集格。
若 $L$ 没有同构于 Fano 拟阵对偶 $F_7^*$ 的子式，则(6.8)、(6.9)描述了 Tutte 复形 $TC(L,a)$ 的所有二维胞腔。［这意味着在表示理论的应用中，上述定义已足够。］一般还需要在某些满足 $\rk(X\wedge Y\wedge Z\wedge T)=r-4$ 的四边形上贴二维胞腔；这种“余秩4型”四边形的定义相当复杂，这里不再描述。
\sourcepage{1834}
\entry{定理6.10}{（同伦定理，Tutte 1958）}复形 $TC(L,a)$ 是1连通的。
其组合含义是：不含 $a$ 的任意两个余点 $X,Y$ 可在 Tutte 图 $TG(L,a)$ 中由不含 $a$ 的路连接；任意两条这样的路相差一系列初等同伦，其形式为 $XYX$、(6.8)中的 $XYZX$、(6.9)中的 $XYZTX$ 或余秩4型。比较定理6.7前的讨论。
此处的表述与 Tutte (1958) 原表述形式不同，内容相同。据 Tutte (1979，第446页)，证明虽长，但本质上是图论和几何性的。除特殊情形外，似乎没有显著简化原证明的方法。其中一种特殊情形是：对所有不含 $a$ 且交为余线的余点 $X,Y$，都有 $\bar X\cup\bar Y\ne L^1-a$。此时由(11.10)(iv)、(11.2)及定理11.14，$L-[a,\hat1]$ 的最高三层构成1连通偏序集；类似定理6.6的证明，应用纤维定理10.5，可将1连通性传给 $TC(L,a)$。寻找一般情形下更简单、更概念化的证明仍有意义。
篇幅不允许详细说明定理6.10在表示理论中的应用，只概要介绍思想。Tutte 正则拟阵刻画的充分性证明对底集大小归纳，故需要删除点 $a$。粗略说，“正则”坐标化定义在余点上，新点 $a$ 处的值沿 Tutte 图中的路从一个余点传到另一个。此时同伦定理保证不同路径不会产生矛盾。定理7.6的证明将更详细展示类似思想。
